% Generated draft: requires mathematical and source review.
\documentclass[11pt]{article}
\usepackage[margin=25mm]{geometry}
\usepackage{fontspec}
\setmainfont{Latin Modern Roman}
\usepackage{amsmath,amssymb,mathtools,unicode-math}
\setmathfont{Latin Modern Math}
\usepackage{enumitem,needspace,xurl,hyperref}
\setlength{\parindent}{0pt}
\setlength{\parskip}{5pt}
\setlength{\emergencystretch}{3em}
\newcommand{\N}{\mathbb{N}}
\newcommand{\Z}{\mathbb{Z}}
\newcommand{\Q}{\mathbb{Q}}
\newcommand{\R}{\mathbb{R}}
\newcommand{\C}{\mathbb{C}}
\newcommand{\F}{\mathbb{F}}
\newcommand{\A}{\mathbb{A}}
\newcommand{\PP}{\mathbb{P}}
\newcommand{\E}{\mathbb{E}}
\newcommand{\eps}{\varepsilon}
\newcommand{\dd}{\,\mathrm{d}}
\begin{document}
\section*{UnsolvedMath ID: 30002686}
\textbf{Source problem code:} catalogue entry
\subsection{Definitions and mathematical statement}
Under the usual Riemannian interpretation of the supplied statement, let \(G\) be a finite-dimensional Lie group, let \(K\leq G\) be a closed subgroup, and let \(M=G/K\) be the corresponding smooth homogeneous space. These conventions explicate the terminology; the source does not explicitly specify the category or impose connectedness, simple connectedness, or effectiveness of the action.

A Riemannian metric \(g\) on \(M\) is \(G\)-invariant if \(L_a^*g=g\) for every \(a\in G\), where \(L_a(bK)=abK\). It is Einstein if there exists a constant \(\lambda\in\R\) such that \(\operatorname{Ric}_g=\lambda g\), and it is nonflat if its Riemann curvature tensor \(\operatorname{Rm}_g\) is not identically zero.

The requested classification, with ``nonflat'' interpreted as a condition on the Einstein metric being sought, is to determine all noncompact homogeneous spaces \(G/K\) for which
\[
\exists\,g\ \exists\,\lambda\in\R:\qquad
 g\text{ is a }G\text{-invariant Riemannian metric},\qquad
 \operatorname{Ric}_g=\lambda g,\qquad
 \operatorname{Rm}_g\not\equiv0.
\]
The source supplies no dimension bound and does not restrict the Einstein constant to be negative. Its background discussion of negative scalar curvature must therefore not be read as an additional hypothesis.

The classification target remains qualified: the source does not specify whether to classify homogeneous presentations \(G/K\), underlying manifolds admitting such presentations, or Einstein metrics, nor whether equivalence means equivariant diffeomorphism, diffeomorphism, isometry, or isometry up to scaling. Moreover, flatness is a property of a metric, so the source's phrase ``nonflat homogeneous spaces'' requires the interpretation stated above or further clarification. These choices require human mathematical review.
\subsection{Short English statement}
Classify noncompact homogeneous spaces that carry an invariant, nonflat Riemannian Einstein metric. The source leaves the equivalence relation and several category conventions unspecified. Its supplied status is \texttt{partially\_solved}; the dated literature assessment concerns negative scalar curvature and is not an independently verified resolution of the full classification problem.
\subsection{Sources}
\begin{itemize}
\item Local UnsolvedMath registry, supplied complete record: permanent ID \(30002686\), title ``Classification of Noncompact Homogeneous Einstein Spaces,'' source code \texttt{catalogue entry}, collection ``Oberwolfach Reports Open Problems.'' The record lists proposed year 2014 and supplies no named proposer.
\item Supplied workshop reference: \emph{Mini-Workshop: Einstein Metrics, Ricci Solitons and Ricci Flow under Symmetry Assumptions} (2015), \emph{Oberwolfach Reports}, pp.\ 2529--2568, DOI \texttt{10.4171/owr/2014/45}. \url{https://doi.org/10.4171/owr/2014/45}. The record describes these reports as documenting mathematical workshops and their open-problem discussions. The supplied year 2015 and the 2014 provenance have not been reconciled here.
\item Supplied literature reference: Bohm--Lafuente, \emph{Non-compact Einstein Manifolds with Symmetry}. \url{https://doi.org/10.1090/jams/1022}.
\item Supplied literature reference: Lauret, \emph{Einstein Solvmanifolds are Standard}. \url{https://annals.math.princeton.edu/2010/172-3/p09}.
\item Literature assessment dated exactly \texttt{2026-08-21}: the supplied curation labels the problem partially solved and asserts that the Alekseevskii conjecture places every negative-scalar-curvature homogeneous Einstein space within the Euclidean-topology and Einstein-solvmanifold framework, while the full classification reduces to the still-incomplete classification of nilsolitons. This is supplied dated curation derived from the OWR record, not an independently verified literature assessment. The record expressly calls for independent verification of time-sensitive status and cited claims.
\item The supplied statement has registry label \texttt{corrected\_verified}, accompanied by a claim of manual semantic, typographic, and LaTeX review confirming a meaning-preserving cleanup. This records the supplied editorial assessment and does not constitute independent mathematical or source verification of this draft.
\end{itemize}
\end{document}
